\documentclass[11pt,a4paper]{article}
\usepackage[T1]{fontenc}
\usepackage{lmodern,amsmath,amssymb,amsthm,mathtools,booktabs,microtype}
\usepackage[margin=26mm]{geometry}
\usepackage[colorlinks=true,linkcolor=blue,citecolor=blue,urlcolor=blue]{hyperref}
\newtheorem{theorem}{Theorem}[section]
\newtheorem{proposition}[theorem]{Proposition}
\newtheorem{lemma}[theorem]{Lemma}
\newtheorem{corollary}[theorem]{Corollary}
\theoremstyle{definition}\newtheorem{definition}[theorem]{Definition}
\newtheorem{remark}[theorem]{Remark}
\newcommand{\Q}{\mathbb Q}\newcommand{\Z}{\mathbb Z}
\newcommand{\vp}{v_p}\newcommand{\rad}{\operatorname{rad}}
\setlength{\emergencystretch}{2em}
\title{Height Bounds from Disjoint Sets of Rational Contacts}
\author{}\date{9 November 2026}
\begin{document}\maketitle
\begin{abstract}
A height bound with exponent $19$ for an odd contact radical at five
rational points implies a bound with exponent $19/r$ for the combined
radical at $5r$ explicitly specified points, for every positive integer
$r$. The proof uses integral linear fractional transformations with
determinant $-1$ and a fixed exceptional set of primes. We give all
linear forms and constants, a compatible rational normalization, and
the resulting inequalities for local congruence arguments. The number
of contacts and the degree of their product increase with $r$; the
height improvement alone does not assert a reduction in the number of
coordinates of a defining polynomial.
\end{abstract}

\section{The height hypothesis and the conclusion}

For $t=m/n\in\Q$ in lowest terms with $n>0$, put
$h(t)=\log\max(|m|,n)$. Valuations are normalized by $v_p(p)=1$.
Let $S_0$ be the set of primes at most $76176$, and put
\[
 H=2^{110}210^{1800}\left(\prod_{p\in S_0}p\right)^{19}+128.
\]
The arithmetic premise used in this paper is the following five-contact
height bound from \cite[Theorem \texttt{quartic-height19}]{Glazer}:
for $s=u/v\in\Q\setminus\{0,1,2401,-343/32\}$ in lowest terms,
\begin{equation}\label{eq:base}
 h(s)\le H M_5(s)^{19},
\end{equation}
where $M_5(s)$ is the product of primes outside $S_0$ for which at
least one of
\[
 v_p(v),\quad v_p(u),\quad v_p(u-v),\quad
 v_p(u-2401v),\quad v_p(32u+343v)
\]
is positive and odd. The contact at infinity is represented by $v$.
All height conclusions below are implications of \eqref{eq:base}.
No stronger height premise is assumed.

For $r\ge1$, define
\[
 \Theta=\left\{0,-\frac12,-1,\frac1{2399},-\frac{32}{407}\right\},
 \qquad \Theta_r=\bigcup_{j=0}^{r-1}(3j+\Theta).
\]
Introduce the following $5r$ primitive integral linear forms:
\begin{align*}
 L_{j,0}(m,n)&=m-3jn,\\
 L_{j,1}(m,n)&=2m+(1-6j)n,\\
 L_{j,2}(m,n)&=m+(1-3j)n,\\
 L_{j,3}(m,n)&=2399m-(7197j+1)n,\\
 L_{j,4}(m,n)&=407m+(32-1221j)n.
\end{align*}
Their zeros $m/n$ are exactly $\Theta_r$.
Write these forms as $L_i=a_i m+b_i n$, using a single index when
the pair $(j,k)$ is immaterial.

\begin{definition}\label{def:exceptional}
The finite set $\Sigma_r$ consists of $S_0$ and all prime divisors of
the nonzero integers
\[
 a_i,\qquad a_i b_{i'}-a_{i'}b_i\ (i\ne i'),\qquad a_i+b_i.
\]
Put
\[
 E_r=\prod_{p\in\Sigma_r\setminus S_0}p,\qquad
 C_r=\log\max(3,9r-10),\qquad H_r=C_r+HE_r^{19}.
\]
For $t\in\Q\setminus\Theta_r$, let $M_r(t)$ be the product of
primes $p\notin\Sigma_r$ for which
$v_p(t-\theta)$ is positive and odd for some $\theta\in\Theta_r$.
\end{definition}

Every integer used to define $\Sigma_r$ is nonzero.
Indeed, the five elements of $\Theta$ are distinct and lie in
$[-1,1/2399]$, so their translates by $3j$ are disjoint.
Also $1\notin\Theta_r$: for $j=0$ this is immediate, and for $j\ge1$
the least element is $3j-1\ge2$. Every $a_i$ is positive.
In particular, all quantities in Definition~\ref{def:exceptional}
are explicitly computable from $r$.

\begin{theorem}\label{thm:amplification}
Assume \eqref{eq:base}. For every positive integer $r$ and every
$t\in\Q\setminus\Theta_r$,
\begin{equation}\label{eq:amplified}
 h(t)\le C_r+HE_r^{19}M_r(t)^{19/r}
       \le H_r M_r(t)^{19/r}.
\end{equation}
Moreover, $M_r(1)=1$.
\end{theorem}

\section{Disjoint radicals and the proof}

\begin{lemma}\label{lem:disjoint}
Let $t=m/n$ be in lowest terms. At a prime $p\notin\Sigma_r$,
at most one of the integers $L_i(m,n)$ is divisible by $p$.
If $p\mid n$, none of them is divisible by $p$.
Consequently $M_r(t)$ is the product of the pairwise coprime
radicals $M_{r,j}(t)$ obtained by using only the five forms with
first index $j$.
\end{lemma}
\begin{proof}
If $p\mid n$, then $m$ and each $a_i$ are units modulo $p$.
Thus $L_i(m,n)$ is a unit. If $p\nmid n$ and two forms vanish
modulo $p$, their determinant vanishes modulo $p$, contrary to
the definition of $\Sigma_r$. Since $a_i$ and $n$ are units in
the latter case,
\[
 v_p\!\left(t+\frac{b_i}{a_i}\right)=v_p(L_i(m,n)).
\]
The assertions about the radicals follow.
\end{proof}

\begin{proof}[Proof of Theorem~\ref{thm:amplification}]
For $0\le j<r$, put
\[
 \phi_j(t)=2+\frac1{t-3j}
          =\frac{2m+(1-6j)n}{m-3jn}.
\]
The numerator and denominator are relatively prime: the matrix
\[
 A_j=\begin{pmatrix}2&1-6j\\1&-3j\end{pmatrix}
 \quad\hbox{has determinant }-1.
\]
Its inverse is
\[
 A_j^{-1}=\begin{pmatrix}3j&1-6j\\1&-2\end{pmatrix}.
\]
The operator norm for the maximum norm is
$\max(3,9j-1)$ for $j\ge0$, interpreting the formula at $j=0$
as $3$. Hence
\begin{equation}\label{eq:height-change}
 h(t)\le h(\phi_j(t))+C_r.
\end{equation}
The five projective contact forms for $\phi_j(t)$ are exactly
$L_{j,0},\ldots,L_{j,4}$, up to signs. In particular, the hypothesis
$t\notin\Theta_r$ excludes every forbidden value and pole of
$\phi_j$. Primes removed by enlarging $S_0$ to $\Sigma_r$ occur
at most once in the radical, regardless of their valuations.
Therefore
\[
 M_5(\phi_j(t))\le E_r M_{r,j}(t).
\]
If $h(t)\le C_r$, the result is immediate. Otherwise
\eqref{eq:base} and \eqref{eq:height-change} give, for every $j$,
\[
 0<h(t)-C_r\le HE_r^{19}M_{r,j}(t)^{19}.
\]
Multiplying these $r$ inequalities and applying
Lemma~\ref{lem:disjoint} proves the first inequality in
\eqref{eq:amplified}. The second follows from $M_r(t)\ge1$.
Finally, every prime dividing a contact value $L_i(1,1)$ belongs
to $\Sigma_r$, so $M_r(1)=1$.
\end{proof}

\begin{remark}
This argument uses the product of disjoint odd-valuation radicals.
It does not replace an odd-valuation radical by an ordinary radical.
At the removed primes, an arbitrary positive valuation contributes
at most one prime factor to the comparison with $M_5$.
\end{remark}

\section{A compatible rational normalization}

The following identities describe how the larger contact set can
enter a local valuation construction. They are algebraic identities;
they do not assert that an arbitrary circuit using them has the
projection property of a defining polynomial.

Let $d=5r$ and set $\epsilon=1$ for odd $r$, $\epsilon=2$ for even
$r$. Thus $d+\epsilon$ is even and $1\le\epsilon\le2$.
Put
\[
 F_r(t)=\prod_{\theta\in\Theta_r}(t-\theta).
\]
Fix a positive rational square $a$, take $t=a\lambda$, and write
$\delta=t-1$. When $\delta F_r(t)\ne0$, define
\begin{equation}\label{eq:normalization}
 R=\frac{F_r(t)}{a^d\delta^{d+\epsilon}},\qquad
 e=\frac{(a\delta)^d}{F_r(t)},\qquad
 q=\frac1{R\delta^4}.
\end{equation}
The corresponding cleared identities are
\begin{equation}\label{eq:cleared}
 F_r(t)e=(a\delta)^d,\qquad q\delta^{4-\epsilon}=e,
 \qquad e\delta^\epsilon=R^{-1}.
\end{equation}

\begin{proposition}\label{prop:normalization}
The following assertions hold.
\begin{enumerate}
\item Every solution of the first identity in \eqref{eq:cleared}
has $F_r(t)\ne0$. If $q=0$ in the first two identities, then
$\delta=0$ and $M_r(t)=1$.
\item If $p\mid M_r(t)$, then $\delta$ is a $p$-adic unit and
$v_p(R)$ is odd. In particular, $v_p(q)=-v_p(R)$.
\item For $\lambda\to+\infty$,
\[
 R\sim a^{-(d+\epsilon)}\lambda^{-\epsilon}>0,
 \qquad q\sim a^{d+\epsilon-4}\lambda^{\epsilon-4}\longrightarrow0.
\]
\item At a fixed prime where $v_p(t)<0$ and is smaller than every
valuation of a nonzero contact, one has
\[
 v_p(q)=d\,v_p(a)-(4-\epsilon)v_p(t).
\]
For fixed $a$, this tends to $+\infty$ as $v_p(\lambda)$ tends
to $-\infty$.
\end{enumerate}
\end{proposition}
\begin{proof}
If $F_r(t)=0$, then $t\in\Theta_r$, whence $\delta\ne0$ and
$(a\delta)^d\ne0$, a contradiction. If $q=0$, the second identity
gives $e=0$, and the first gives $\delta=0$; Theorem~\ref{thm:amplification}
gives $M_r(t)=1$.

At a consumed prime, exactly one contact difference has positive odd
valuation, all other contact differences are units, and $t-1$ is a
unit. These facts follow from the determinant and $L_i(1,1)$
conditions in $\Sigma_r$. Since $v_p(a)$ is even, the valuation
$v_p(F_r(t))-d\,v_p(a)$ is odd. This proves the second assertion.
The third follows by comparing the leading terms in
\eqref{eq:normalization}. For the fourth, every contact difference
and $\delta$ has valuation $v_p(t)$. Thus $v_p(e)=d\,v_p(a)$,
and the second identity in \eqref{eq:cleared} proves the formula.
\end{proof}

\begin{proposition}\label{prop:prime-pattern}
Suppose $\lambda=u/v$ with $v$ a positive integral square and
$u\in\Z$, and suppose the $d$ numbers
\[
 p_\theta=u-\frac{\theta}{a}v\qquad(\theta\in\Theta_r)
\]
are positive primes. Assume they avoid $\Sigma_r$ and every prime
dividing the numerator or denominator of $a$. If they are distinct,
then $R$ is positive for $\lambda$ sufficiently large and its
reduced denominator is a square. Its primes of odd valuation
are exactly the $p_\theta$, with valuation one. If in addition each
$p_\theta$ is a square in $\Q_2$, $\Q_3$, and $\Q_{17}$, then
so is $R$.
\end{proposition}
\begin{proof}
Direct substitution gives
\[
 R=\frac{\prod_{\theta\in\Theta_r}p_\theta}
          {v^d\delta^{d+\epsilon}}.
\]
The denominator is a rational square. We check that it is a unit at
each $p_\theta$. Such a prime cannot divide $v$: otherwise it also
divides $u$ and hence all the distinct primes $p_{\theta'}$, since
all contact and scale denominators are units there. Thus $a\lambda$
reduces to $\theta$ at this prime, and $\delta$ reduces to
$\theta-1\ne0$. Its denominator is therefore a unit, and the
valuation of $R$ is one. At every other prime the valuation is even.
This proves the assertion about the reduced denominator and the odd
valuations. Positivity follows from $d+\epsilon$ even and the
positive numerator. The final assertion follows because a product
of local squares divided by a nonzero square is a square.
\end{proof}

The prime-pattern hypothesis in Proposition~\ref{prop:prime-pattern}
is the same finite-pattern requirement used in the arithmetic
construction: the contacts and scales are fixed first, rational
coefficients are cleared by a fixed square denominator, and a
finite set of primes is avoided \cite{Sun,Arithmetic}.
Increasing $r$ changes this finite list and the associated constants.
An explicit numerical cutoff must consequently be recalculated.

\section{The divisibility exponents required by height arguments}

\begin{lemma}\label{lem:two-scales}
Let $\alpha>0$, $H_*>0$, and suppose a radical $M$ satisfies
$h(s)\le H_*M(s)^\alpha$ at two nonzero rational values
$s=\lambda$ and $s=T^2\lambda$, where $T\ne0$.
If $N\ne0$ is an integer and
\[
 M(\lambda)^k\mid N,\qquad M(T^2\lambda)^k\mid N,
\]
then
\[
 h(T)\le H_*|N|^{\alpha/k}.
\]
The two radicals need not be relatively prime.
\end{lemma}
\begin{proof}
The elementary height inequalities give
\[
 2h(T)=h(T^2)\le h(T^2\lambda)+h(\lambda).
\]
Each radical is at most $|N|^{1/k}$, proving the result.
\end{proof}

\begin{corollary}\label{cor:residual}
Assume \eqref{eq:base}. Suppose $B\ge1$, $A>0$, $J\ge0$,
$T,\lambda\in\Q^\times$, and
\[
 B^2\le A h(T)+J.
\]
Assume both $\lambda$ and $T^2\lambda$ avoid $\Theta_r$.
Let an integer residual $\Delta$ satisfy
\[
 |\Delta|\le5B^4,\qquad
 M_r(\lambda)^k\mid\Delta,\qquad
 M_r(T^2\lambda)^k\mid\Delta.
\]
If $rk>38$, then all solutions with $B$ sufficiently large in terms
of $r,k,A,J,H_r$ satisfy $\Delta=0$.
\end{corollary}
\begin{proof}
If $\Delta\ne0$, Theorem~\ref{thm:amplification} and
Lemma~\ref{lem:two-scales} imply
\[
 B^2\le AH_r 5^{19/(rk)}B^{76/(rk)}+J.
\]
Since $76/(rk)<2$, this bounds $B$ above by a constant depending
only on the displayed data.
\end{proof}

\begin{corollary}\label{cor:input}
Under the same hypotheses, suppose an integer $N\ne0$ instead
satisfies
\[
 |N|\le2\exp(h(\beta))B,\qquad
 M_r(\lambda)^\mu\mid N,\qquad
 M_r(T^2\lambda)^\mu\mid N.
\]
If $2r\mu>19$, then
\[
 \log B\le \frac{19}{2r\mu-19}\,h(\beta)+C,
\]
where $C$ depends only on $r,\mu,A,J,H_r$.
\end{corollary}
\begin{proof}
Put $\gamma=19/(r\mu)<2$. The preceding argument gives
\[
 B^2\le AH_r2^\gamma\exp(\gamma h(\beta))B^\gamma+J.
\]
If $B^2\le2J$, the assertion follows by increasing $C$.
Otherwise absorb $J$ into $B^2/2$, take logarithms, and divide by
$2-\gamma$. The coefficient of $h(\beta)$ becomes
$\gamma/(2-\gamma)=19/(2r\mu-19)$.
\end{proof}

The least positive integers meeting these two strict inequalities are
\[
 k_r=\left\lfloor\frac{38}{r}\right\rfloor+1,
 \qquad \mu_r=\left\lfloor\frac{19}{2r}\right\rfloor+1.
\]
Some useful values are:
\begin{center}
\begin{tabular}{rrrrrrrrrrrr}\toprule
$r$&1&2&3&4&5&6&7&10&13&20&39\\\midrule
$5r$&5&10&15&20&25&30&35&50&65&100&195\\
$k_r$&39&20&13&10&8&7&6&4&3&2&1\\
$\mu_r$&10&5&4&3&2&2&2&1&1&1&1\\\bottomrule
\end{tabular}
\end{center}
These are requirements on divisibility of integer residuals. A formal
approximation formula must separately be proved accurate enough to
produce those divisibilities. Lower exponents can shorten such a
formula, but the enlarged product $F_r$ requires additional arithmetic
operations. Neither this table nor Theorem~\ref{thm:amplification}
is a coordinate count for a polynomial definition of the integers.

\section{Exact finite verification}

The accompanying standard-library program
\nolinkurl{check_contact_height_amplification.py} reconstructs the
linear forms for the values of $r$ in the table, verifies their
primitivity and pairwise distinct zeros, and factors every integer
defining $\Sigma_r$. It records the exceptional primes larger than
$76176$, the exact product $E_r$, and the height-transformation norm.
It also checks the integer threshold inequalities and the normalized
valuation identities. This calculation verifies the finite algebra;
the arithmetic premise \eqref{eq:base} remains explicit.

\begin{thebibliography}{9}
\bibitem{Glazer} E. Glazer and contributors,
\emph{Diophantine definition of the integers in the rationals},
\texttt{manuscripts/quartic\_v9.tex}, commit
\texttt{14e8f601fe0d9173d4a5f084aedc3e07066a7688},
\url{https://github.com/elliotglazer/diophantine-z-q}.
\bibitem{Sun} Z.-W. Sun,
\emph{86 unknowns for Hilbert's tenth problem over the rationals},
arXiv:2610.20691v2, Theorem 4.1,
\url{https://arxiv.org/abs/2610.20691v2}.
\bibitem{Arithmetic} J. Brox and contributors,
\emph{Arithmetic dependencies and the corrected source argument},
\texttt{arithmetic\_dependencies\_20261108\_v8.tex},
Sections 2--3,
\url{https://github.com/JoseBrox/diophantine-z-q_improvements}.
\end{thebibliography}
\end{document}
